\Needspace{8\baselineskip}
\Article{203}
\textsc{Aufgabe.} \textit{Sind die Formen \(\Phi,\varphi\) äquivalent, so sind alle Transformationen der einen in die andere anzugeben.}

\textit{Lösung.} Sind diese Formen nur auf eine Art äquivalent (\textit{d.h.}, entweder nur eigentlich oder nur uneigentlich), so suchen wir nach Artikel~196 eine Transformation von \(\varphi\) in \(\Phi\), etwa \(\alpha,\beta,\gamma,\delta\). Offenbar kann es keine anderen als dazu gleichartige Transformationen geben. Sind aber \(\varphi,\Phi\) sowohl eigentlich als auch uneigentlich äquivalent, so suchen wir zwei ungleichartige Transformationen, \textit{d.h.} eine eigentliche und eine uneigentliche, etwa \(\alpha,\beta,\gamma,\delta\) und \(\alpha',\beta',\gamma',\delta'\). Jede andere Transformation ist mit der einen oder der anderen gleichartig. Sei also die Form \(\varphi\) gleich \((a,b,c)\), ihre Determinante \(=D\) und der größte gemeinsame Teiler von \(a,2b,c\) wie stets zuvor \(m\); ferner sollen \(t,u\) alle Zahlen durchlaufen, die \(tt-Duu=mm\) erfüllen. Dann sind im ersten Fall alle Transformationen von \(\varphi\) in \(\Phi\) in der ersten der folgenden Formeln, I, enthalten, im zweiten Fall entweder in der ersten, I, oder in der zweiten, II.
\[
\begin{array}{@{}lll@{}}
\mathrm{I}\;\ldotp\,\ldotp\,\ldotp\,\ldotp\,\ldotp\,\ldotp\,\ldotp&\frac1m(\alpha t-(\alpha b+\gamma c)u),&\frac1m(\beta t-(\beta b+\delta c)u)\\[1ex]
&\frac1m(\gamma t+(\alpha a+\gamma b)u),&\frac1m(\delta t+(\beta a+\delta b)u)\\[1.5ex]
\mathrm{II}\;\ldotp\,\ldotp\,\ldotp\,\ldotp\,\ldotp\,\ldotp\,\ldotp&\frac1m(\alpha't-(\alpha'b+\gamma'c)u),&\frac1m(\beta't-(\beta'b+\delta'c)u)\\[1ex]
&\frac1m(\gamma't+(\alpha'a+\gamma'b)u),&\frac1m(\delta't+(\beta'a+\delta'b)u)
\end{array}
\]
\SourceLeaf{202}
\textit{Beispiel.} - Gesucht sind alle Transformationen von $(129,92,65)$ in $(42,59,81)$. In Artikel~195 haben wir gefunden, dass diese nur uneigentlich äquivalent sind, und im nächsten Artikel die uneigentliche Transformation $-47,-56,73,87$ der ersten in die zweite ermittelt. Daher werden alle Transformationen von $(129,92,65)$ in $(42,59,81)$ durch die Formel
\[
-(47t+421u),\quad -(56t+503u).\quad 73t+653u,\quad 87t+780u
\]
dargestellt, wobei $t$, $u$ alle Zahlen durchlaufen, die $tt-79uu=1$ erfüllen. Diese werden dargestellt durch
\[
\begin{aligned}
\pm t&=\tfrac12\bigl((80+9\sqrt{79})^e+(80-9\sqrt{79})^e\bigr)\\
\pm u&=\tfrac{1}{2\sqrt{79}}\bigl((80+9\sqrt{79})^e-(80-9\sqrt{79})^e\bigr)
\end{aligned}
\]
wobei für $e$ alle nichtnegativen ganzen Zahlen zu nehmen sind.

\Needspace{8\baselineskip}
\Article{204}
Offenbar wird die allgemeine Formel, die alle Transformationen darstellt, umso \textit{einfacher}, je einfacher die Ausgangstransformation ist, aus der wir sie ableiten. Da wir frei wählen können, von welcher Transformation wir ausgehen, lässt sich die allgemeine Formel oft einfacher machen: Wir leiten aus der zuerst gefundenen Formel eine einfachere Transformation ab, indem wir $t$, $u$ bestimmte Werte geben, und bilden daraus eine andere Formel. Setzen wir \textit{z.B.} in der Formel aus dem Beispiel des vorhergehenden Artikels $t=80$, $u=-9$, so erhalten wir eine einfachere Transformation als die, von der wir ausgegangen waren, nämlich $29,47,-37,-60$. Daraus folgt die allgemeine Formel $29t-263u$, $47t-424u$, $-37t+337u$, $-60t+543u$. Haben wir also nach den vorhergehenden Regeln eine allgemeine Formel gefunden, können wir versuchen, ob sich, wenn wir $t$, $u$ die bestimmten Werte $\pm t',\pm u';\ \pm t'',\pm u''$ usw. geben, eine einfachere Transformation ergibt als die, aus der die Formel abgeleitet wurde. In diesem Fall lässt sich aus jener Transformation eine einfachere Formel herleiten. --- Bei der Beurteilung der Einfachheit bleibt allerdings ein gewisser Spielraum. Wenn es die Mühe lohnte, könnten wir dafür eine feste Regel aufstellen und in der Folge $t',u';\ t'',u''$ usw. auch \textit{Grenzen} angeben, jenseits derer die Transformationen ständig weniger einfach werden. Dann brauchten wir nicht weiterzugehen, sondern nur innerhalb dieser Grenzen zu suchen. Da aber nach unseren Methoden die einfachste Transformation meist entweder so\SourceLeaf{203}fort oder indem wir $t$, $u$ die Werte $\pm t',\pm u'$ geben, auftritt, lassen wir diese Untersuchung der Kürze wegen weg.

\Needspace{8\baselineskip}
\Article{205}
\textsc{Aufgabe.} \textit{Alle Darstellungen einer gegebenen Zahl $M$ durch eine gegebene Formel $axx+2bxy+cyy$ mit positiver nichtquadratischer Determinante $=D$ finden.}

\textit{Lösung.} Zunächst bemerken wir, dass sich die Suche nach Darstellungen durch nicht teilerfremde Werte von $x$, $y$ genauso wie oben für Formen mit negativer Determinante (Artikel~181) auf den Fall zurückführen lässt, in dem teilerfremde Werte der Unbestimmten gesucht werden. Es wäre daher überflüssig, dies hier zu wiederholen. Damit Darstellungen durch teilerfremde Werte von $x,y$ möglich sind, muss $D$ ein quadratischer Rest von $M$ sein. Sind alle Werte des Ausdrucks $\sqrt D\;(\mathrm{mod.}\,M)$ gleich $N,-N,N',-N',N'',-N''$ usw. (die wir so wählen dürfen, dass keiner $>\tfrac12M$ ist), so gehört jede Darstellung von $M$ durch die gegebene Form zu einem dieser Werte. Wir müssen also zuerst diese Werte ermitteln und dann die zu jedem einzelnen gehörigen Darstellungen untersuchen. Darstellungen zum Wert $N$ gibt es nur, wenn die Formen $(a,b,c)$ und $(M,N,\frac{NN-D}{M})$ eigentlich äquivalent sind. Ist das der Fall, so suchen wir eine beliebige eigentliche Transformation der ersten in die zweite, etwa $\alpha,\beta,\gamma,\delta$. Dann haben wir eine Darstellung der Zahl $M$ durch $(a,b,c)$, die zum Wert $N$ gehört, nämlich $x=\alpha$, $y=\gamma$, und alle zu diesem Wert gehörigen Darstellungen werden durch die Formel
\[
x=\tfrac1m\bigl(\alpha t-(\alpha b+\gamma c)u\bigr),\qquad
y=\tfrac1m\bigl(\gamma t+(\alpha a+\gamma b)u\bigr)
\]
dargestellt. Dabei bezeichnet $m$ den größten gemeinsamen Teiler von $a,2b,c$, und $t$, $u$ durchlaufen alle Zahlen, die $tt-Duu=mm$ erfüllen. --- Offenbar wird diese allgemeine Formel umso einfacher, je einfacher die Transformation $\alpha,\beta,\gamma,\delta$ ist, aus der sie abgeleitet wurde. Deshalb ist es sinnvoll, vorher nach dem letzten Artikel die einfachste Transformation von $(a,b,c)$ in $(M,N,\frac{NN-D}{M})$ zu ermitteln und daraus die Formel abzuleiten. --- Auf genau dieselbe Weise lassen sich die zu den übrigen Werten $-N,N',-N'$ usw. gehörigen Darstellungen, falls es welche gibt, durch allgemeine Formeln angeben.

\textit{Beispiel.} Gesucht sind alle Darstellungen von $585$ durch die Formel
\SourceLeaf{204}
$42xx+62xy+21yy$. Bei Darstellungen durch nicht teilerfremde Werte von $x$, $y$ ist sofort klar, dass nur solche möglich sind, bei denen $x$, $y$ den größten gemeinsamen Teiler $3$ haben, denn $585$ ist nur durch das eine Quadrat $9$ teilbar. Haben wir also alle Darstellungen von $\frac{585}{9}$, \textit{d.h.} $65$, durch die Form $42x'x'+62x'y'+21y'y'$ gefunden, in denen $x'$ und $y'$ teilerfremd sind, so erhalten wir daraus alle Darstellungen von $585$ durch $42xx+62xy+21yy$, in denen $x$ und $y$ nicht teilerfremd sind, indem wir $x=3x'$, $y=3y'$ setzen. Die Werte von $\sqrt{79}\;(\mathrm{mod.}\,65)$ sind $\pm12,\pm27$. Als Darstellung von $65$ zum Wert $-12$ finden wir $x'=2$, $y'=-1$. Daher werden alle Darstellungen von $65$ zu diesem Wert durch $x'=2t-41u$, $y'=-t+53u$ und alle \textit{daraus} entstehenden Darstellungen von $585$ durch $x=6t-123u$, $y=-3t+159u$ gegeben. Ebenso finden wir die allgemeine Formel für alle Darstellungen von $65$ zum Wert $+12$, nämlich $x'=22t-199u$, $y'=-23t+211u$, und die Formel für alle daraus entstehenden Darstellungen von $585$, nämlich $x=66t-597u$, $y'=-69t+633u$. Zu den Werten $+27$ und $-27$ gehört aber keine Darstellung von $65$. --- Um die Darstellungen von $585$ durch teilerfremde Werte von $x$, $y$ zu finden, müssen wir zuerst die Werte von $\sqrt{79}\;(\mathrm{mod.}\,585)$ ermitteln; sie sind $\pm77,\pm103,\pm157,\pm248$. Zu den Werten $\pm77,\pm103,\pm248$ finden wir keine Darstellung. Zum Wert $-157$ gehört aber die Darstellung $x=3$, $y=1$, aus der wir die allgemeine Formel für alle zu diesem Wert gehörigen Darstellungen ableiten: $x=3t-114u$, $y=t+157u$. Ebenso finden wir als Darstellung zum Wert $+157$ die Werte $x=83$, $y=-87$ und als Formel, die alle gleichartigen enthält, $x=83t-746u$, $y=-87t+789u$. Wir haben also vier allgemeine Formeln, unter denen sämtliche Darstellungen von $585$ durch $42xx+62xy+21yy$ enthalten sind:
\[
\begin{aligned}
x&=6t-123u&\qquad y&=-3t+159u\\
x&=66t-597u& y&=-69t+633u\\
x&=3t-114u& y&=t+157u\\
x&=83t-746u& y&=-87t+789u
\end{aligned}
\]
wobei $t,u$ alle ganzen Zahlen durchlaufen, die $tt-79uu=1$ erfüllen.

Bei besonderen Anwendungen der vorhergehenden Untersuchungen über Formen mit positiver nichtquadratischer Deter\SourceLeaf{205}minante halten wir uns der Kürze wegen nicht auf. Jeder kann sie ohne Schwierigkeit selbst durchführen, ähnlich wie in den Artikeln~176, 182. Wir wenden uns gleich den Formen mit positiver quadratischer Determinante zu, die allein noch übrig sind.
